% ==============================================================================
% v3 BDI §3 — The cumulative carry $P_a$ as analytical backbone
% Standalone fragment; to be \input{} into v3 main.tex.
% ==============================================================================
%
% Macros the v3 preamble must provide (inherited from v2 main.tex unless flagged NEW):
%
%   \Kp, \NT, \Sing, \Chain, \MB, \TM, \Cmid, \Csing, \KK, \Z, \Zn,
%   \g, \eps, \labr, \labl                                            (v2)
%
%   \newcommand{\Pcum}{P^{\mathrm{cum}}}                              (NEW)
%   \newcommand{\Pn}{\mathcal{P}}                                     (NEW; chain polytope)
%   \newcommand{\hw}{\mathrm{hw}}                                     (NEW)
%   \newcommand{\piNT}{\pi^{\NT}}                                     (NEW)
%   \newcommand{\Rec}{\mathcal{R}}                                    (NEW; recording set)
%   \newcommand{\HW}{\mathrm{HW}}                                     (NEW)
%   \DeclareMathOperator{\wt}{wt}                                     (NEW)
%
% Theorem environments are those declared in v2: theorem, proposition, lemma,
% corollary, definition, example, remark.  Numbering continues from \section{}.
% ==============================================================================

\section{The cumulative carry $P_a$ as analytical backbone}\label{sec:carry}

%-------------------------------------------------------------------------------
\subsection{Roadmap}\label{ssec:carry-roadmap}

Throughout this section $\g$ is of type $B_n$, $n \ge 2$, and $\pi \in \Kp(\infty)$
has the chain--sing coordinates
\[
   \bigl((M_a, B_a, T_a)_{a=1}^{n-1},\; S,\; \piNT\bigr)
\]
of \S\ref{sec:prelim}. We isolate a single scalar, the \emph{cumulative carry}
\[
   P_a \;:=\; \sum_{b=1}^{a} 2(B_b - T_b), \qquad P_0 := 0,
\]
and show that it carries the entire short-simple analytic content of the
chain--singleton decomposition at $B_n$. The carry plays five structurally
distinct roles, the first four forced by the single left-to-right bracket scan
of Theorem~\ref{thm:carry-A}, the fifth obtained from the first by Farkas-polar:

\begin{enumerate}[label=(R\arabic*),leftmargin=2.5em]
  \item \textbf{Descent recording} (\S\ref{ssec:carry-descent}). $P_a$ is one-sided
    monotone along forward $\widetilde{e}_n$-descent. This obstructs any local
    witness for the image $\Rec(\pi^{\hw})$, closing v3 OPEN-3 by impossibility.
  \item \textbf{Singleton cross-chain coupling} (\S\ref{ssec:carry-cross}). The
    unique cross-chain condition in the $B_n$-highest indicator is the carry
    sum lift $[\,S \le P_{n-1}\,]$, type-uniform in $n$.
  \item \textbf{Carry-recursive factorisation} (\S\ref{ssec:carry-factor}). Per-chain
    MB and TM tests couple chain $a$ to all prior chains through $P_{a-1}$;
    the $B_n$-HW indicator factors carry-recursively, not as a strict tensor.
  \item \textbf{Polytope completeness} (\S\ref{ssec:carry-polytope}). The chain-side
    HW polytope $\Pn_n$ has exactly $2n - 3$ non-redundant carry-derived facets,
    namely $\{L_a, U_a : 2 \le a \le n-1\} \cup \{E\}$, with $L_1$ degenerate
    and $U_1$ redundant.
  \item \textbf{Weight-space simplicial cone} (\S\ref{ssec:weight-simplicial}).
    Under the chain-to-weight map $\Phi$, the image cone $\Kp_n \subset
    \mathbb{R}^n$ is simplicial with exactly $n$ facets; of the $2n - 3$
    chain-side carry-derived facets, only $E$ survives projection, and the
    chain-to-weight contraction has ratio $(2n - 3)/n \to 2$.
\end{enumerate}

External polyhedral and algebraic shadows are recorded in
\S\ref{ssec:kobayashi-fences} (Kobayashi) and \S\ref{ssec:carry-shadows}
(Watanabe, Meereboer, Frohmader); the closing remark (\S\ref{ssec:carry-closing})
collects the five roles as five views of one object.

%-------------------------------------------------------------------------------
\subsection{The cumulative carry $P_a$ and Theorem A}\label{ssec:carry-Pa}

We recall the chain--sing alphabet of \S\ref{sec:prelim}. For each
$a = 1, \dots, n-1$, the chain-$a$ multiplicities $(M_a, B_a, T_a)$ count the
roots $(E_a,\, E_a - E_n,\, E_a + E_n)$ respectively; the singleton multiplicity
$S$ counts $E_n$. The non-touching block $\piNT$ contributes no bracket symbols
at $\alpha_n$ and is fixed under the $B_n$-action; it plays no role in this
section, and we drop it from notation after Lemma~\ref{lem:carry-mono}.

The CST bracket sequence $S_n(\pi)$ for the short simple $\alpha_n$ is the
left-to-right concatenation, for $a = 1, \dots, n-1$, of the per-chain block
\[
   \labr{E_a}^{M_a}\;\labl{E_a - E_n}^{2 B_a}\;\labr{E_a + E_n}^{2 T_a}\;\labl{E_a}^{M_a},
\]
followed by the singleton block $\labr{E_n}^{S}$. Standard cancellation defines
$\eps_n(\pi)$.

\begin{definition}[Cumulative carry]\label{def:carry}
Set $P_0 := 0$ and, for $a = 1, \dots, n-1$,
\[
   P_a \;:=\; P_{a-1} + 2(B_a - T_a) \;=\; \sum_{b=1}^{a} 2(B_b - T_b).
\]
We occasionally write $\Pcum_a$ when the cumulative nature is the focus.
\end{definition}

Up to sign, $P_a = -\langle \alpha_n^\vee,\, \wt(\pi)|_{\le a}\rangle$: the carry
records twice the cumulative deficit of $E_n$-content over chains $1, \dots, a$.

\begin{lemma}[One-sided monotonicity]\label{lem:carry-mono}
Each forward chain-factor descent step ($\widetilde{e}_n$ acting on the rightmost
surviving $\labr{}$ in $S_n(\pi)$) adjusts $P_a$ as follows:
\begin{itemize}[leftmargin=2em]
  \item $\Sing$: $S \mapsto S - 1$; no change to any $P_a$.
  \item $\MB(a)$: $(M_a, B_a, T_a) \mapsto (M_a - 1, B_a + 1, T_a)$; $P_b \mapsto P_b + 2$ for $b \ge a$.
  \item $\TM(a)$: $(M_a, B_a, T_a) \mapsto (M_a + 1, B_a, T_a - 1)$; $P_b \mapsto P_b + 2$ for $b \ge a$.
\end{itemize}
No letter decreases any $P_a$ under forward descent.
\end{lemma}

\begin{proof}
Direct from Definitions~\ref{def:step-types} and~\ref{def:carry}.
\end{proof}

The next statement is Theorem~A in carry form, with $P_a$ as the load-bearing
object. It is the master lemma on which \S\ref{ssec:carry-descent}--\S\ref{ssec:weight-simplicial} rest.

\begin{theorem}[Theorem~A in carry form]\label{thm:carry-A}
$\eps_n(\pi) = 0$ if and only if
\begin{align}
   (\HW_a) \quad & M_a \le P_{a-1} \;\text{ and }\; M_a \le P_a,
   && a = 1, \dots, n-1, \label{eq:carry-HWa}\\
   (\HW_{\mathrm{sing}}) \quad & S \le P_{n-1}. \label{eq:carry-HWsing}
\end{align}
\end{theorem}

\begin{proof}
A left-to-right scan of $S_n(\pi)$ tracking the running open-paren count $Q$. By
induction on $a$, on entry to chain $a$ we have $Q = P_{a-1}$ provided no
$\labr{}$ has survived; the chain-$a$ subscan survives a $\labr{}$ in exactly
the cases excluded by~\eqref{eq:carry-HWa}, after rewriting the second of the
two original inequalities ($M_a + 2 T_a \le P_{a-1} + 2 B_a$) as $M_a \le P_a$.
The singleton block then processes against $Q = P_{n-1}$, giving
\eqref{eq:carry-HWsing}.
\end{proof}

We use the clean form~\eqref{eq:carry-HWa} throughout the remainder of the section.

%-------------------------------------------------------------------------------
\subsection{Descent role: no local witness}\label{ssec:carry-descent}

Theorem~B (\S\ref{sec:phaseB}) makes the descent step
$\pi \mapsto (\pi^{\hw}, R(\pi))$ a bijection
\[
   \Kp(\infty)/B_n \;\longrightarrow\; \bigsqcup_{\pi^{\hw}} \{\pi^{\hw}\} \times \Rec(\pi^{\hw}),
\]
where $\Rec(\pi^{\hw})$ is the set of valid recording words. The natural question
(v3 OPEN-3) is whether $\Rec(\pi^{\hw})$ admits a \emph{local} characterisation:
a finite list of prefix/suffix inequalities on a candidate word $R$ that
decides membership without replaying the descent.

\begin{proposition}[No local witness]\label{prop:no-local-witness}
There is no signed prefix or suffix invariant on the recording word $R$ whose
non-negativity (or vanishing) characterises $R \in \Rec(\pi^{\hw})$.
\end{proposition}

\begin{proof}
By Lemma~\ref{lem:carry-mono}, every forward step adds $0$ or $+2$ to each $P_b$
and never decreases any. The reverse-played word therefore monotonically
\emph{drains} $P_b$. For a candidate word $R$ that is not in $\Rec(\pi^{\hw})$,
intermediate states during reverse-play can drive some $P_b$ arbitrarily
negative with no later letter able to compensate; the failure cannot be diagnosed
by any signed counter on $R$ alone. The only correct test is the oracle: $R \in
\Rec(\pi^{\hw})$ iff reverse-playing $R$ from $\pi^{\hw}$ stays coordinate-nonnegative
and forward descent of the resulting Kostant partition recovers $R$.
\end{proof}

Proposition~\ref{prop:no-local-witness} is supported by a three-probe falsification
arc on the candidate inequalities (i)~prefix carry $\ge 0$, (i$'$)~pair-aware
prefix carry, (iii)~per-chain suffix balance. All three give zero false positives
and monotonically growing false negatives ($268, 360, 410$) as the candidate
tightens: the signature of an impossibility, not a tuning miss.

\begin{remark}[Asymmetric mirror]\label{rmk:asymmetric-mirror}
Lemma~\ref{lem:carry-mono} is the analytical heart of an
\emph{asymmetric-mirror} phenomenon: forward and
reverse descent are not exchanged by a local involution because the alphabet
$\{\Sing, \MB(a), \TM(a)\}$ has no signed compensator at the carry level. The
chain-factor alphabet is rank-3, but the carry sees only a $\{0, +2\}$ bit. The
signed slack data $(t_0^{(i)}, r^{(i)})$ that Azenhas's type-AII recording
exploits has no chain-factor analogue: $\MB(a)$ and $\TM(a)$ are distinct
chain-factor bits but contribute identically to $P_a$. This is the
rank-3-chain-factor-vs-rank-1-strip distinction of \S\ref{ssec:scope} at the
analytic level, and it is why no local witness exists for $\Rec(\pi^{\hw})$.
\end{remark}

%-------------------------------------------------------------------------------
\subsection{Projection role: Theorem E}\label{ssec:carry-cross}

The second role of the carry is to mediate the unique cross-chain condition in
the $B_n$-highest indicator.

\begin{theorem}[Singleton cross-chain coupling]\label{thm:carry-E}
Let $\pi \in \Kp(\infty)$ have chain--sing coordinates
$((M_a, B_a, T_a)_{a=1}^{n-1}, S)$. Then $\pi$ is $B_n$-highest if and only if
\[
   (\HW_a) \quad M_a \le P_{a-1}\ \text{and}\ M_a \le P_a, \quad a = 1, \dots, n-1,
\]
\[
   (\HW_{\mathrm{sing}}) \quad S \le P_{n-1} \;=\; \sum_{a=1}^{n-1} 2(B_a - T_a).
\]
In particular, the only cross-chain condition is the cumulative carry bound
\[
   \mathrm{Cross}(M, B, T, S) \;=\; [\,S \le P_{n-1}\,] \;=\; \Bigl[\, S \le \sum_{a=1}^{n-1} 2(B_a - T_a) \,\Bigr],
\]
which depends only on the final cumulative carry and is type-uniform in $n$.
\end{theorem}

\begin{proof}
The biconditional is Theorem~\ref{thm:carry-A}. The proof there is a single
left-to-right scan whose recurrence and per-block surviving-bracket count have
no $n$-dependence beyond the chain index, and whose singleton block always
processes against $P_{n-1}$.
\end{proof}

\begin{remark}[The NT block contributes zero]\label{rmk:NT-drop}
The non-touching block $\piNT$ contributes no bracket symbols at $\alpha_n$
(every non-touching root has $\langle \cdot, \alpha_n^\vee\rangle = 0$ and trivial
$\alpha_n$-string), so it appears in neither $S_n(\pi)$ nor any $P_a$. The cross-chain
sum~\eqref{eq:carry-HWsing} therefore reads against the chain--sing coordinates
alone, and $\piNT$ is free under the $B_n$-action. This is the silent assumption
underlying Theorem~\ref{thm:carry-E} and Corollary~\ref{cor:carry-factor}: every analytic
identity at $\alpha_n$ is invariant under arbitrary changes to $\piNT$.
\end{remark}

\begin{example}[Singleton draws on cumulative carry]\label{ex:carry-cross-B4}
At $B_4$ the configuration $(M, B, T, S) = ((0,0,0), (0,1,0), (0,0,0),\, 2)$ has
$P_1 = 0$, $P_2 = 2$, $P_3 = 2$, so $S = 2 = P_{n-1}$ and the configuration is
$B_4$-highest with all $(\HW_a)$ trivially satisfied. The singleton capacity
is supplied by chain~$2$ alone, confirming that the cross-chain condition is
genuinely cumulative, not pointwise per-chain.
\end{example}

Empirically, on the $1716$ chain+sing configurations at $B_3$, content $\le 6$,
the sum-lift $[S \le P_{n-1}]$ is the unique candidate (among pointwise
$\bigwedge_a [S \le 2(B_a - T_a)]$, sum lift, cumulative-and, min, max) with zero
false negatives; the same dataset falsifies the strict tensor factorisation
``$\chi_1(M_1, B_1, T_1) \cdot \chi_2(M_2, B_2, T_2) \cdot \mathrm{Cross}$'' with
$34$ inconsistent $(\mathrm{chain}_1, \Delta P_2, S)$-keys. A $B_4$ spot check
($3003$ configurations at content $\le 5$) gives $3003/3003$. Full write-up:
\url{proofs/2026-05-20-pnu-Bn-cross-chain.md}.

%-------------------------------------------------------------------------------
\subsection{Carry-recursive factorisation}\label{ssec:carry-factor}

Theorem~\ref{thm:carry-E} admits an immediate corollary in factorised form.

\begin{corollary}[Carry-recursive factorisation]\label{cor:carry-factor}
The $B_n$-HW indicator factors as
\[
   [\,B_n\text{-HW}(M, B, T, S)\,]
   \;=\; \prod_{a=1}^{n-1} \chi_a(M_a, B_a, T_a;\, P_{a-1}) \,\cdot\, [\,S \le P_{n-1}\,],
\]
with $\chi_a = [M_a \le P_{a-1}] \wedge [M_a \le P_a]$ and the carry update
$P_a = P_{a-1} + 2(B_a - T_a)$, $P_0 = 0$.
\end{corollary}

The factorisation is \emph{carry-recursive}, not strict tensor: $\chi_a$ depends
on the incoming scalar $P_{a-1}$ as well as on $(M_a, B_a, T_a)$, and the
singleton factor uses $P_{n-1}$. As in Remark~\ref{rmk:NT-drop}, the NT block
contributes neither to the per-chain factors $\chi_a$ nor to the singleton
bound, so it is absent from the corollary. Corollary~\ref{cor:carry-factor} is
the unified replacement for the descent-recording and projection-cross-chain
results of the old v3 outline: both follow from a single left-to-right bracket
scan acting on a single scalar.

%-------------------------------------------------------------------------------
\subsection{Polytope completeness (Theorem~F)}\label{ssec:carry-polytope}

The fourth role of the carry is structural. The system $(\HW_a) \cup
(\HW_{\mathrm{sing}})$ of Theorem~\ref{thm:carry-A} is a list of $2n - 1$ linear
inequalities; we ask how many of them are non-redundant facets of the
chain-side HW polytope, and which.

\begin{definition}[Chain HW polytope]\label{def:chain-polytope}
The \emph{chain HW polytope} at $B_n$ is
\[
   \Pn_n \;:=\; \bigl\{ (M, B, T, S) \in \Z_{\ge 0}^{3(n-1) + 1} : \pi \text{ is } B_n\text{-highest} \bigr\}
\]
(viewing $\Pn_n$ as a lattice polytope, equivalently the intersection of
non-negativity with the carry system).
\end{definition}

Label the $2n - 1$ carry-derived inequalities of $\Pn_n$ as
\[
   L_a : M_a \le P_{a-1}, \quad U_a : M_a \le P_a \qquad (1 \le a \le n - 1),
   \qquad E : S \le P_{n-1}.
\]

\begin{theorem}[Polytope completeness]\label{thm:carry-F}
Let $n \ge 2$. The chain HW polytope $\Pn_n$ has exactly $2n - 3$ non-redundant
carry-derived facets, namely
\[
   \{L_a, U_a : 2 \le a \le n-1\} \;\cup\; \{E\}.
\]
The boundary inequalities at $a = 1$ are special:
\begin{enumerate}[label=(\roman*),leftmargin=2em]
  \item $L_1$ is \emph{degenerate}: it is non-redundant as an inequality but is
    saturated on every $\pi \in \Pn_n$, forcing $M_1 = 0$ identically. Equivalently,
    $\Pn_n$ lies in the affine hyperplane $\{M_1 = 0\}$.
  \item $U_1$ is \emph{redundant}: it is implied by $L_1$ together with non-negativity
    of $M_2$ (for $n \ge 3$) or of $S$ (for $n = 2$), via $L_2 + (M_2 \ge 0)$
    respectively $E + (S \ge 0)$.
\end{enumerate}
\end{theorem}

\begin{proof}
\emph{(i)~$L_1$ degenerate.} $L_1$ reads $M_1 \le P_0 = 0$. With $M_1 \ge 0$,
this forces $M_1 = 0$ on $\Pn_n$. Non-redundancy: the configuration
$(M_1, B_1, T_1) = (1, 1, 0)$ (other coordinates zero) has $P_1 = 2$ and
satisfies every $L_a, U_a$ for $a \ge 2$ and $E$ (each reads against $M_a = 0$
or $S = 0$ with $P_{n-1} = 2$), but violates $L_1$. So $L_1$ is needed; on
$\Pn_n$ it is always tight.

\emph{(ii)~$U_1$ redundant.} By (i), $M_1 = 0$ on $\Pn_n$, so $U_1$ becomes the
slack inequality $0 \le P_1$. If $n = 2$, fence $E$ reads $S \le P_1$ and
non-negativity gives $P_1 \ge S \ge 0$. If $n \ge 3$, fence $L_2$ reads
$M_2 \le P_1$ and non-negativity gives $P_1 \ge M_2 \ge 0$. Either way, the
other inequalities together with non-negativity force $U_1$.

\emph{(iii)~$\{L_a, U_a : 2 \le a \le n-1\} \cup \{E\}$ are non-redundant.} For
each, we exhibit a chain configuration that satisfies every \emph{other}
inequality of the system but violates the named one (existence of such a witness
proves non-redundancy; saturation then has a relative-interior point, giving a
codimension-$1$ face).

\smallskip
\noindent\textit{Witness for $L_a$ ($2 \le a \le n-1$):} set $M_a = 1$,
$B_a = 1$, all other coordinates zero. Then $P_b = 0$ for $b < a$ and $P_b = 2$
for $b \ge a$. Every $L_b, U_b$ for $b < a$ has $M_b = 0 \le 0$; $U_a$ reads
$1 \le 2$; every $L_b, U_b$ for $b > a$ reads $0 \le 2$; $E$ reads $0 \le 2$;
$L_a$ reads $1 \le 0$ and is the unique violation.

\smallskip
\noindent\textit{Witness for $U_a$ ($2 \le a \le n-1$):} set $M_a = 1$,
$B_{a-1} = 1$, $T_a = 1$, all other coordinates zero. Then $P_b = 0$ for
$b < a-1$, $P_{a-1} = 2$, $P_b = 0$ for $b \ge a$. $L_a$ reads $1 \le 2$;
$U_a$ reads $1 \le 0$ and is the unique violation; all others are slack.

\smallskip
\noindent\textit{Witness for $E$:} set $S = 1$, all chain coordinates zero. Then
$P_{n-1} = 0$; every $L_a, U_a$ reads $0 \le 0$; $E$ reads $1 \le 0$ and is the
unique violation.

\smallskip
\noindent To upgrade non-redundancy to codimension-$1$, type-uniform
relative-interior witnesses are constructed in
\texttt{verify\_facet\_interiors.py} (verified at $n = 3, \dots, 7$); we sketch:
for $L_a$, take $B_b = 1$ ($1 \le b \le a-1$), $M_a = 2(a-1)$, $B_a = 1$, giving
$L_a$ saturated and all others strict; for $U_a$ ($a \ge 3$), take the same
prefix with $M_a = 2(a-2)$, $T_a = 1$, $B_a = 0$, giving $U_a$ saturated and all
others strict; for $U_2$, take $B_1 = 2$, $M_2 = 2$, $T_2 = 1$; for $E$, take
$B_b = 1$ ($1 \le b \le n-1$), $S = 2(n-1)$.
\end{proof}

\begin{remark}[Boundary behaviour]\label{rmk:boundary-asymmetry}
The asymmetry at $a = 1$ is structural rather than incidental: the carry
recursion starts at $P_0 = 0$, so $L_1$ is the strongest possible mid-bound
($M_1 \le 0$) and $U_1$ is the corresponding slack form ($P_1 \ge 0$) that the
next downstream constraint enforces. Internally ($2 \le a \le n-1$), $P_{a-1}$
and $P_a$ both range over positive values and neither is universally tighter,
so both $L_a$ and $U_a$ contribute honest facets.
\end{remark}

\begin{remark}[Computational verification]\label{rmk:carry-F-verification}
The script \texttt{verify\_witnesses.py} verifies the three-class structure
($L_1$ degenerate, $U_1$ redundant, all other carry-derived inequalities
non-redundant) for $n \le 7$, and \texttt{classify\_fences.py} enumerates all
chain configurations of total content $\le C$ at $B_n$, $n \le 6$, classifying
each inequality by saturation count, strict-satisfaction count, and explicit
non-redundancy count. The polytope-completeness data table (Day-$28$ proof,
\url{proofs/2026-05-21-bdi-kobayashi-faces.md}) records the explicit
non-redundancy counts at $(n, C) = (2, 8), (3, 6), (4, 6), (5, 4), (6, 3)$.
\end{remark}

The polytope role completes the chain-side unification: $P_a$ is the analytical
backbone in \S\ref{ssec:carry-Pa}, the descent potential in
\S\ref{ssec:carry-descent}, the cross-chain coupling in
\S\ref{ssec:carry-cross}, the recursive scalar in \S\ref{ssec:carry-factor},
and the facet-counting object in \S\ref{ssec:carry-polytope}. The same scan
that proves Theorem~\ref{thm:carry-A} proves all four. The fifth role
(\S\ref{ssec:weight-simplicial}) is obtained from the same chain-side data by
Farkas-polar under the chain-to-weight projection $\Phi$.

%-------------------------------------------------------------------------------
\subsubsection{Weight-space image is simplicial (Theorem~G)}\label{ssec:weight-simplicial}

Theorem~\ref{thm:carry-F} counts the carry-derived facets of $\Pn_n$ on the
chain side, in the $3(n-1) + 1$ chain--singleton coordinates. The
representation-theoretic object of interest is, however, the image of $\Pn_n$
under the chain-to-weight map
\[
   \Phi : \Pn_n \times \Z_{\ge 0}^{\NT} \;\longrightarrow\; X^+(B_n) \otimes \mathbb{R}
   \;\cong\; \mathbb{R}^n,
\]
which records the highest weight $\nu = \sum_\beta m_\beta\,\beta$ in the
$(E_1, \dots, E_n)$ basis. Concretely $\lambda_a = M_a + B_a + T_a +
(\text{NT contribs at } E_a)$ for $a \le n - 1$, and
$\lambda_n = S + \sum_{a=1}^{n-1}(T_a - B_a)$. Let
$\Kp_n := \Phi(\Pn_n \times \Z_{\ge 0}^{\NT}) \subset \mathbb{R}^n$.

\begin{theorem}[Weight-space image polytope]\label{thm:carry-G}
Let $n \ge 2$. The image cone $\Kp_n$ is the rational polyhedral cone in
$\mathbb{R}^n$ cut out by exactly $n$ facets:
\begin{equation}
\begin{aligned}
   &\lambda_1 + \lambda_2 + \cdots + \lambda_k \;\ge\; 0
   && (1 \le k \le n - 2),\\
   &\sum_{i=1}^{n} \lambda_i \;\ge\; 0, \qquad
   \lambda_n \;\le\; \sum_{i=1}^{n-1} \lambda_i.
\end{aligned}\label{eq:carry-G}
\end{equation}
The $n$ inequalities are pairwise non-redundant. In particular $\Kp_n$ is a
\emph{simplicial cone} in $\mathbb{R}^n$ (the minimum possible facet count for a
full-dimensional polyhedral cone).
\end{theorem}

Theorem~\ref{thm:carry-G} should be compared with Theorem~\ref{thm:carry-F}:
the chain side carries $2n - 3$ non-redundant carry-derived facets, the
weight side carries only $n$. The chain-to-weight projection is contractive of
ratio $(2n - 3)/n \to 2$ as $n \to \infty$: roughly half the chain
combinatorial structure is invisible from weight space. Of the $2n - 3$ chain
facets $\{L_a, U_a : 2 \le a \le n-1\} \cup \{E\}$, only $E$ survives
projection (as $\lambda_n \le \sum_{i \le n-1} \lambda_i$); every $L_a$ and
$U_a$ collapses, and the NT block contributes \emph{zero} new facets while
\emph{absorbing} the partial-sum facet at $k = n - 1$.

\begin{proof}[Proof sketch]
\emph{(i) Polar duality framing.} To find the facet normals of an image cone
$\Phi(C)$, one polars: $c \in \Kp_n^*$ iff $\langle c, \Phi(x)\rangle \le 0$
for all $x$ in the source cone, iff $\Phi^T c \in \Pn_n^*$, iff $\Phi^T c$ is a
non-negative combination of the rows of the source H-representation (Farkas).
Facets of $\Kp_n$ are then read off as extreme rays of $\Kp_n^*$.

\emph{(ii) Explicit verification at $B_3$.} The chain-side H-representation has
rows $L_2, U_2, E$ together with coordinate non-negativities. Writing
$c = (c_1, c_2, c_3)$ and computing $\Phi^T c$ component-by-component in the
order $(B_1, T_1, M_2, B_2, T_2, S, N^-_{12}, N^+_{12})$ gives an eight-equation
system in the multipliers $(p, q, r) = (\mu_{L_2}, \mu_{U_2}, \mu_E)$ and the
non-negativity multipliers. The NT rows force $c_1 \le c_2$ and
$c_1 + c_2 \le 0$ directly; the chain rows reduce, after eliminating $(p, q)$,
to the bound $|c_3| \le -c_2$. The polar cone is
\[
   \Kp_3^* \;=\; \{c \in \mathbb{R}^3 : c_1 \le c_2 \le c_3,\;\; c_2 + c_3 \le 0\},
\]
whose three extreme rays $(-1, 0, 0), (-1, -1, -1), (-1, -1, 1)$ give exactly
the three image facets $\lambda_1 \ge 0$, $\sum_i \lambda_i \ge 0$, and
$\lambda_3 \le \lambda_1 + \lambda_2$.

\emph{(iii) Type-uniform mechanism.} The same polar argument carries over for
$n \ge 2$. The NT pairs $(a, b)$ with $a < b \le n - 1$ contribute
$c_a \le c_b$ and $c_a + c_b \le 0$, producing the total order
$c_1 \le c_2 \le \cdots \le c_{n-1} \le 0$. The chain rows $L_a, U_a$ pull
back, via $\Phi^T$, into the interior of the polar relative to the dominant
binding constraints from $E$ and from non-negativity: their multipliers can
always be set to zero, so $L_a, U_a$ are Farkas-redundant in the projected
H-rep (this is the formal sense in which $\Phi(L_a)$ and $\Phi(U_a)$ cover the
full image cone). Only $\mu_E$ remains as a non-trivial extreme contribution,
yielding the final binding inequality $c_{n-1} + c_n \le 0$. The polar cone
$\Kp_n^* = \{c : c_1 \le \cdots \le c_n,\; c_{n-1} + c_n \le 0\}$ has exactly
$n$ extreme rays $\rho_k = -(e_1 + \cdots + e_k)$ for $1 \le k \le n - 2$,
$\rho_{\mathrm{sum}} = -(e_1 + \cdots + e_n)$, and $\rho_E = -(e_1 + \cdots +
e_{n-1}) + e_n$, giving the $n$ facets of~\eqref{eq:carry-G}. NT absorption of
the would-be facet $\sum_{i \le n-1} \lambda_i \ge 0$ is automatic: $\piNT$
contributes negative-coordinate moves into $\lambda_{n-1}$ (and into earlier
$\lambda_a$), and the partial sum at $k = n - 1$ is also implied by the
combination of the full sum and the $E$-facet.\footnote{Verified
computationally at $B_2, B_3, B_4, B_5$; see
\url{proofs/2026-05-21-weight-space-projection.md}.}
\end{proof}

\begin{remark}[Weight-indistinguishability of MB and TM]\label{rmk:weight-indistinguishability}
Theorem~\ref{thm:carry-G} provides the abstract justification for the
MB/TM weight-indistinguishability used in Theorem~B (Phase~B): on weight
space, the chain inequalities $L_a$ and $U_a$ produce the same shadow,
because both pull back via $\Phi^T$ into the Farkas interior of the polar
cone. The signed difference $L_a - U_a = 2(T_a - B_a)$ is precisely the
$\lambda_n$-content of chain $a$; this differential is visible in $\lambda_n$
but does not generate an additional facet, so the $\{L_a, U_a\}$ pair is
weight-indistinguishable in the polytope sense. Whatever finer information
the chain framework records is by construction invisible from the highest-weight
polytope alone, and must be recovered via the carry recursion of
Corollary~\ref{cor:carry-factor}.
\end{remark}

%-------------------------------------------------------------------------------
\subsection{External shadow: Kobayashi's BDI fences}\label{ssec:kobayashi-fences}

Kobayashi~\cite{Kobayashi2604} establishes that for orthogonal symmetric pairs
of BDI type, in particular $(O(n+1), O(n))$, branching multiplicities are
governed by piecewise-linear ``fence'' hyperplanes of the form
$\xi_i \pm \delta\nu_j = \pm 1/2$ in infinitesimal-character space, decomposing
the parameter region into convex chambers on which the multiplicity function is
locally polynomial. The fences exist abstractly, but Kobayashi does not give
an explicit list in highest-weight coordinates. Theorem~\ref{thm:carry-F}
(chain side, $2n - 3$ facets) together with Theorem~\ref{thm:carry-G}
(image side, $n$ facets) supply the first explicit determination of the BDI
fences in the chain--sing coordinates of \S\ref{sec:prelim}: the $n$ facets
of $\Kp_n$ in~\eqref{eq:carry-G} are precisely the stability fences of
$(O(n+1), O(n))$-branching, and the additional $n - 3$ chain-side carry walls
$\{L_a, U_a : 2 \le a \le n - 1\}$ describe the sub-fence structure that
becomes invisible on projection. Read this way, the carry framework is the
\emph{combinatorial content} of Kobayashi's analytic stability statement: the
chain side records fine multiplicity-level data, the weight side records the
coarse fence chamber structure, and the chain-to-weight contraction is the
explicit form of Kobayashi's piecewise-linearity. The precise coordinate
match between Kobayashi's $\xi_i + \delta\nu_j = \pm 1/2$ form and the
chain-side wall list $\{L_a, U_a, E\}$ is expected by representation theory but
not yet verified term by term; we record this as the open question OQ-KFENCE
and defer it to v4.

%-------------------------------------------------------------------------------
\subsection{Other external shadows}\label{ssec:carry-shadows}

Three further external papers give shadows that the carry framework illuminates.
(Kobayashi's BDI fences are treated separately in \S\ref{ssec:kobayashi-fences}.)

\paragraph{Algebraic existence.} Watanabe \cite{Watanabe2407} (Prop.~4.3.2, \S5)
establishes existence of the abstract orthogonal projection
$p_\nu : V(\nu) \to V^\imath_{\mathrm{BDI}}(\nu)$ under Bao--Wang preferred
parameters, type-uniformly in $n$. Theorem~\ref{thm:carry-E} gives the explicit
shape of $p_\nu$ in the chain-factor basis: a carry-recursive product whose only
cross-chain ingredient is the cumulative bound $[S \le P_{n-1}]$. The shadow
says $p_\nu$ exists; the carry says what it looks like.

\paragraph{Base case.} Meereboer \cite{Meereboer2510} (Lemma~6.2, \S6) treats the
$n = 1$ base case from a combinatorial leading-term recipe. At $n = 1$,
Theorem~\ref{thm:carry-E} collapses to $S \le P_0 = 0$, i.e., $S = 0$, matching
Meereboer's recipe. Meereboer's Lemmas~4.3, 4.5, 7.4 collapse at multiplicity
$> 1$; the higher-multiplicity, higher-rank carry-recursive structure fills the
gap above $n = 1$.

\paragraph{An alternate combinatorial model.} Frohmader \cite{Frohmader2312v2}
gives a different combinatorial model for BDI branching as a sum, over partitions
$\mu$ with even rows, of modified Littlewood--Richardson coefficients with an
$O_n$-highlighted-box flag. The decomposition is not of chain-factor type:
Frohmader's $\mu$-index aggregates the chain data and the singleton in a way
that has no direct counterpart in $\{L_a, U_a, E\}$. Numerical multiplicities
agree by representation theory, so a small-rank numerical match would not
be informative; a structural correspondence at the level of indices would be,
and is open. We record the construction as an alternate combinatorial viewpoint;
it is \emph{not} an independent confirmation of the chain-factor framework. The
carry-derived facet count $2n - 3$ is specific to the chain coordinates of
\S\ref{sec:prelim}.

%-------------------------------------------------------------------------------
\subsection{Closing remark}\label{ssec:carry-closing}

The cumulative carry $P_a$ is the unified analytical object at BDI.
Definition~\ref{def:carry} introduces it; Lemma~\ref{lem:carry-mono} records its
one-sided monotonicity; Theorem~\ref{thm:carry-A} characterises $B_n$-highest
weight in terms of it; the four chain-side roles of
\S\ref{ssec:carry-descent}--\S\ref{ssec:carry-polytope} are mechanically forced
by the same left-to-right bracket scan; the weight-space simplicial role of
\S\ref{ssec:weight-simplicial} is the Farkas-polar of the chain-side picture
under $\Phi$; the external shadows of \S\ref{ssec:kobayashi-fences} and
\S\ref{ssec:carry-shadows} agree.

Three points warrant explicit notice. \emph{First}, the unification is
structural, not cosmetic: the same scan proves
Theorem~\ref{thm:carry-A}, Proposition~\ref{prop:no-local-witness},
Theorem~\ref{thm:carry-E}, Corollary~\ref{cor:carry-factor}, and
Theorem~\ref{thm:carry-F}, with Theorem~\ref{thm:carry-G} obtained from the
chain-side H-representation by Farkas-polar under $\Phi$; there is no shorter
path. \emph{Second}, the carry is type-uniform: the recurrence $P_a = P_{a-1}
+ 2(B_a - T_a)$ has no $n$-dependence beyond the chain index, the singleton
bound always reads against $P_{n-1}$, the chain-side facet count $2n - 3$ has
the same shape at every $n \ge 2$, and the image-cone facet count $n$ matches
the simplicial minimum. \emph{Third}, the framework's scope coincides with the
v2 chain-factor decomposition's scope, the $|a_{ij}| \le 2$ regime of the
Day-$21$ trichotomy (\S\ref{ssec:scope}); the $G_2$ obstruction at $|a_{ij}|
\ge 3$ has multiple interior bracket positions and breaks the carry recursion.

The carry was, in retrospect, available in the bracket-scan argument of
\S\ref{sec:carry} from the start. The work of v3 is to recognise that the same
scalar plays five distinct roles, and to compress the exposition so that each
role is a view of the same object.

% ==============================================================================
% Placeholder bibitems for v3.bib integration (Rick: integrate later).
% ==============================================================================
%
% \bibitem{Watanabe2407}     H. Watanabe.  arXiv:2407.07280v3.
% \bibitem{Meereboer2510}    M. Meereboer.  arXiv:2510.17655.
% \bibitem{Kobayashi2604}    T. Kobayashi.  arXiv:2604.22262.
% \bibitem{Frohmader2312v2}  A. Frohmader.  arXiv:2312.11295v2.
%
% (Other internal citations -- Watanabe2107, Watanabe2110, etc. -- live in v2.bib.)
